%% ============================================================
%% Teil IV: Das Projekt – vollständige Implementierung
%% Band 2 – Kapitel 15
%% ============================================================

\part{Das Projekt}

% ─────────────────────────────────────────────────────────────
\chapter{Wir bauen robotik-uno-q als ROS2-Paket –
  vom Sketch bis zum autonomen Fahren}
% ─────────────────────────────────────────────────────────────

\begin{lernziele}
Du baust \texttt{stufe2\_ros2/} zu einem vollständigen ROS2-Paket um,
implementierst alle drei Nodes in ihrer finalen Form,
verknüpfst sie über die Launch-Datei
und führst einen vollständigen Systemtest durch.
Am Ende fährt der Roboter autonom.
\end{lernziele}

\section{Situation}

In den vorangehenden Kapiteln sind entstanden:
\begin{itemize}
\item \texttt{arduino\_bridge\_node}: liest Serial, published \texttt{sensor\_msgs/Range}
\item \texttt{navigation\_node}: subscribed Sensor-Topics, published \texttt{/cmd\_vel}
\item \texttt{motor\_node}: subscribed \texttt{/cmd\_vel}, sendet Serial-Befehle
\item \texttt{roboter.launch.py}: startet alle drei
\end{itemize}

Jetzt fehlt noch das ROS2-Paket-Gerüst:
\texttt{package.xml}, \texttt{setup.py} und die Eintragung
der Nodes als Executables.

\begin{aufgabeBox}
Forme \texttt{stufe2\_ros2/} in ein vollständiges ROS2-Python-Paket um,
installiere es mit \texttt{colcon build}
und starte den kompletten Stack mit \texttt{ros2 launch}.
\end{aufgabeBox}

\section{ROS2-Paketstruktur}

\begin{lstlisting}[language=bash]
stufe2_ros2/                          # wird umbenannt zu robotik_uno_q/
├── package.xml                       # Paket-Metadaten
├── setup.py                          # Python-Paket + Executables
├── setup.cfg
├── resource/
│   └── robotik_uno_q                 # Marker-Datei (leer)
├── robotik_uno_q/
│   ├── __init__.py
│   ├── arduino_bridge_node.py
│   ├── navigation_node.py
│   ├── motor_node.py
│   └── sensor_daten.py
├── config/
│   └── sensor_config.yaml
└── launch/
    └── roboter.launch.py
\end{lstlisting}

\section{package.xml}

\begin{lstlisting}[language=XML]
<?xml version="1.0"?>
<?xml-model href="http://download.ros.org/schema/package_format3.xsd"?>
<package format="3">
  <name>robotik_uno_q</name>
  <version>0.2.0</version>
  <description>
    ROS2-Paket für robotik-uno-q: HC-SR04-Sensoren, L298N-Motorsteuerung.
  </description>
  <maintainer email="w.becker@schule.de">Wolfgang Becker</maintainer>
  <license>MIT</license>

  <depend>rclpy</depend>
  <depend>sensor_msgs</depend>
  <depend>geometry_msgs</depend>
  <depend>std_srvs</depend>

  <exec_depend>python3-pyserial</exec_depend>
  <exec_depend>python3-yaml</exec_depend>
  <exec_depend>python3-numpy</exec_depend>

  <export>
    <build_type>ament_python</build_type>
  </export>
</package>
\end{lstlisting}

\section{setup.py}

\begin{lstlisting}
from setuptools import setup, find_packages
from glob import glob

package_name = 'robotik_uno_q'

setup(
    name=package_name,
    version='0.2.0',
    packages=find_packages(exclude=['test']),
    data_files=[
        ('share/ament_index/resource_index/packages',
            ['resource/' + package_name]),
        ('share/' + package_name,           ['package.xml']),
        ('share/' + package_name + '/config', glob('config/*.yaml')),
        ('share/' + package_name + '/launch', glob('launch/*.py')),
    ],
    install_requires=['setuptools'],
    zip_safe=True,
    maintainer='Wolfgang Becker',
    maintainer_email='w.becker@schule.de',
    description='HC-SR04 + L298N als ROS2-Paket',
    license='MIT',
    entry_points={
        'console_scripts': [
            'arduino_bridge_node = robotik_uno_q.arduino_bridge_node:main',
            'navigation_node     = robotik_uno_q.navigation_node:main',
            'motor_node          = robotik_uno_q.motor_node:main',
        ],
    },
)
\end{lstlisting}

\section{Der vollständige arduino\_bridge\_node}

\begin{lstlisting}
#!/usr/bin/env python3
# robotik_uno_q/arduino_bridge_node.py -- Version 1.0
"""
ArduinoBridgeNode: Liest HC-SR04-Daten vom Arduino UNO Q
über Serial und published sie als sensor_msgs/Range.
"""

from __future__ import annotations
import re
import time
import serial
import rclpy
from rclpy.node import Node
from rclpy.qos import QoSProfile, ReliabilityPolicy, HistoryPolicy
from sensor_msgs.msg import Range
from .sensor_daten import SensorDaten, SensorPuffer

SENSOR_QOS = QoSProfile(
    reliability=ReliabilityPolicy.BEST_EFFORT,
    history=HistoryPolicy.KEEP_LAST,
    depth=1
)

class ArduinoBridgeNode(Node):
    """Liest Serial-Daten vom Arduino und published Range-Nachrichten."""

    _MUSTER = re.compile(r'R=(\d+)\s+cm\s+L=(\d+)\s+cm')

    def __init__(self) -> None:
        super().__init__('arduino_bridge')

        # Parameter
        self.declare_parameter('port',           '/dev/arduino')
        self.declare_parameter('baudrate',       9600)
        self.declare_parameter('puffer_groesse', 5)

        self._port      = self.get_parameter('port').value
        self._baudrate  = self.get_parameter('baudrate').value

        # Publisher
        self._pub_r = self.create_publisher(Range, '/sensor/rechts', SENSOR_QOS)
        self._pub_l = self.create_publisher(Range, '/sensor/links',  SENSOR_QOS)

        # Puffer
        puffer = self.get_parameter('puffer_groesse').value
        self._puffer = SensorPuffer(fenstergroesse=puffer)

        # Serial-Verbindung
        self._ser: serial.Serial | None = None
        self._verbinde()

        # Timer: 10 Hz
        self._timer = self.create_timer(0.1, self._lese_und_publiziere)
        self.get_logger().info(
            f"ArduinoBridgeNode v1.0 bereit. Port={self._port}")

    def _verbinde(self) -> None:
        try:
            self._ser = serial.Serial(
                self._port, self._baudrate, timeout=1.0)
            time.sleep(2.0)   # Arduino-Neustart abwarten
            self.get_logger().info(f"Verbunden: {self._port}")
        except serial.SerialException as e:
            self.get_logger().error(f"Serial-Fehler: {e}")
            self._ser = None

    def _lese_und_publiziere(self) -> None:
        if self._ser is None or not self._ser.is_open:
            self._verbinde()
            return

        try:
            zeile = self._ser.readline().decode('utf-8', errors='replace')
            daten = self._parse(zeile)
            if daten is None:
                return
            self._puffer.hinzufuegen(daten.rechts_cm, daten.links_cm)
            if not self._puffer.ist_bereit():
                return
            r_m, l_m = self._puffer.gefilterter_wert()
            stamp = self.get_clock().now().to_msg()
            self._pub_r.publish(self._range_msg(r_m, 'sensor_rechts', stamp))
            self._pub_l.publish(self._range_msg(l_m, 'sensor_links',  stamp))

        except serial.SerialException as e:
            self.get_logger().warning(f"Serial-Unterbrechung: {e}")
            self._ser = None

    def _parse(self, zeile: str) -> SensorDaten | None:
        m = self._MUSTER.match(zeile.strip())
        if m:
            return SensorDaten(
                rechts_cm=int(m.group(1)),
                links_cm=int(m.group(2)))
        return None

    @staticmethod
    def _range_msg(meter: float, frame: str, stamp) -> Range:
        msg = Range()
        msg.header.stamp    = stamp
        msg.header.frame_id = frame
        msg.radiation_type  = Range.ULTRASOUND
        msg.field_of_view   = 0.26
        msg.min_range       = 0.02
        msg.max_range       = 4.00
        msg.range           = float(meter) / 100.0
        return msg

    def destroy_node(self) -> None:
        if self._ser and self._ser.is_open:
            self._ser.close()
        super().destroy_node()

def main(args: list[str] | None = None) -> None:
    rclpy.init(args=args)
    node = ArduinoBridgeNode()
    try:
        rclpy.spin(node)
    except KeyboardInterrupt:
        pass
    finally:
        node.destroy_node()
        rclpy.shutdown()

if __name__ == '__main__':
    main()
\end{lstlisting}

\section{Build und Systemtest}

\begin{loesungBox}
\begin{lstlisting}[language=bash]
# 1. ROS2-Workspace vorbereiten
cd ~/ros2_ws                # oder: mkdir ~/ros2_ws && cd $_
mkdir -p src
# robotik-uno-q/stufe2_ros2 als Symlink oder Kopie:
ln -s ~/robotik-uno-q/stufe2_ros2 src/robotik_uno_q

# 2. Paket bauen
source /opt/ros/jazzy/setup.bash
colcon build --packages-select robotik_uno_q
source install/setup.bash

# 3. Stack starten (Arduino angesteckt!)
ros2 launch robotik_uno_q roboter.launch.py

# 4. Prüfen
ros2 node list
# /arduino_bridge  /navigation  /motor

ros2 topic hz /sensor/rechts
# average rate: 10.000

ros2 topic echo /sensor/rechts --once

ros2 param list /arduino_bridge
\end{lstlisting}
\end{loesungBox}

\begin{projektBox}[robotik-uno-q -- Projektstand nach Band~2]
\begin{tabular}{ll}
\textbf{stufe1\_mcu/} & Arduino-Sketches (unverändert aus Band~1) \\
\textbf{stufe2\_ros2/} & Vollständiges ROS2-Paket \texttt{robotik\_uno\_q} \\
\quad arduino\_bridge\_node & liest HC-SR04, published Range \\
\quad navigation\_node & entscheidet, published /cmd\_vel \\
\quad motor\_node & sendet Serial, stellt /motor/stopp bereit \\
\quad roboter.launch.py & startet alles mit einem Befehl \\
\textbf{Roboter} & fährt autonom, weicht Hindernissen aus \\
\end{tabular}

\bigskip
\textbf{Nächster Band:} C++ für Robotik und ROS2 -- \texttt{motor\_node}
wird als C++-Lifecycle-Node neu geschrieben für bessere Echtzeiteigenschaften.
\end{projektBox}

\begin{merksatzBox}
\textbf{Merksatz:}
\texttt{package.xml} deklariert Abhängigkeiten.
\texttt{setup.py entry\_points} registriert Nodes als ausführbare Befehle.
\texttt{colcon build} kompiliert und installiert.
\texttt{source install/setup.bash} macht die neuen Nodes für \texttt{ros2 run} sichtbar.
\end{merksatzBox}

\begin{robotikBox}
Das Paket \texttt{robotik\_uno\_q} folgt exakt der ROS2-Konvention.
Es kann auf GitHub veröffentlicht,
in einem anderen Workspace geklont
und mit \texttt{rosdep install} automatisch mit Abhängigkeiten versehen werden.

In Band~3 kommt \texttt{motor\_node} als C++-Node dazu:
Python und C++ laufen im selben Stack,
kommunizieren über \texttt{/cmd\_vel} --
ohne dass einer von dem anderen weiß.
Das ist ROS2.
\end{robotikBox}
